\insertImage{6b670e2e-7b15-43a0-be2a-03c5853f4a58.png}{Aquarelle d'un dessin d'une balance, pesant les avantages et les inconvénients.}

\chapter{Les critiques du modèle d'entreprise libérée}

Un conseil que je donne à tous les convaincus du modèle — et je m'inclus dans le lot : écoutez vos critiques les plus intelligents. Pas les ricaneurs de fond de salle, mais ceux qui ont de vrais arguments. Ils voient les failles avant vous, et ils vous les offrent gratuitement.

Ce chapitre leur donne la parole, sérieusement. Trois familles d'arguments : économiques, sociaux, et de gouvernance. Je vous préviens : certains vont piquer. C'est le but.

\section{Les arguments économiques et financiers}

Le premier argument des sceptiques est le plus simple : et si la liberté coûtait plus cher qu'elle ne rapporte ?

Leur raisonnement mérite mieux qu'un haussement d'épaules. Sans hiérarchie ni contrôles, disent-ils, la discipline s'érode : la coordination se dégrade, les délais glissent, les erreurs coûtent. La recherche sur les organisations autogérées confirme d'ailleurs que la coordination sans hiérarchie a un coût réel, en temps et en énergie, et que ce coût grimpe avec la taille\footnote{Lee, M. Y., \& Edmondson, A. C. (2017). Self-managing organizations: Exploring the limits of less-hierarchical organizing. Research in Organizational Behavior, 37, 35-58.}. Même Valve, la vitrine du sans-chef, a été analysée sous cet angle : son modèle fonctionne, mais dans des conditions particulières — talents d'exception, marges confortables, métier créatif — qui ne se transposent pas partout\footnote{Puranam, P., \& Håkonsson, D. D. (2015). Valve's way. Journal of Organization Design, 4(2), 2-4.}.

Deuxième volet économique : le pilotage. Sans objectifs imposés ni reporting classique, comment planifier, mesurer, rendre des comptes ? La question n'est pas rhétorique quand vos interlocuteurs s'appellent investisseurs, banquiers ou commissaires aux comptes — des gens dont le métier exige des informations précises, comparables, opposables. Une entreprise libérée doit produire cette lisibilité par d'autres moyens, et certaines n'y arrivent jamais.

Troisième volet, la viabilité à long terme : des coûts de formation et d'accompagnement durablement plus élevés, et le risque qu'une entreprise concentrée sur l'épanouissement relâche sa tension commerciale\footnote{Stø, T., \& Ims, K. J. (2018). Freedom Inc. revisited: creating a freedom-based and healthier organization. Society and Business Review.}. On a vu au chapitre sur la performance que des contre-exemples existent. Mais des contre-exemples ne sont pas une démonstration — et sur ce point, reconnaissons que le débat empirique n'est pas tranché.

\section{Les arguments sociaux et culturels}

La deuxième famille de critiques est celle qui me touche le plus, parce qu'elle défend des gens que le discours de la libération oublie.

Premier argument : l'autonomie ne convient pas à tout le monde, et ce n'est pas un défaut. Des personnes préfèrent un cadre structuré, des attentes explicites, des responsabilités bornées — par tempérament, par étape de vie, ou simplement par choix\footnote{Hsieh, T. (2010). Delivering happiness: A path to profits, passion, and purpose. Hachette Books.}. Les cultures aussi diffèrent : le rapport à la hiérarchie varie énormément d'un pays à l'autre, et un modèle né quelque part ne se greffe pas partout\footnote{Hofstede, G. (2001). Culture's Consequences: Comparing Values, Behaviors, Institutions and Organizations Across Nations. Sage.}. Un modèle qui ne fonctionne qu'avec un certain type de personnes doit au minimum se demander ce qu'il fait des autres.

Deuxième argument, le plus grave à mes yeux : la liberté peut cacher une intensification du travail. C'est la critique du surmenage, celle-là même qui m'a poussé à écrire ce livre. Quand l'entreprise ne fixe plus de limites, c'est au salarié de fixer les siennes — et beaucoup n'y arrivent pas, surtout dans une culture d'enthousiasme obligatoire où lever le pied ressemble à une trahison. Le contrôle n'a pas disparu : il s'est intériorisé. Un chef, on peut lui dire non. Comment dit-on non à une « culture » ?

Troisième argument : l'informel favorise les dominants. Sans rôles définis, qui prend l'ascendant ? Les plus à l'aise à l'oral, les mieux connectés, les plus disponibles — c'est-à-dire, statistiquement, toujours les mêmes profils. La hiérarchie formelle a au moins un mérite : elle se voit, donc elle se conteste. Le pouvoir informel, lui, n'offre aucune prise. Une entreprise libérée qui ne se dote pas de garde-fous explicites contre ce phénomène reproduit les inégalités qu'elle prétendait abolir — en pire, parce qu'invisibles\footnote{Laloux, F. (2014). Reinventing organizations: A guide to creating organizations inspired by the next stage of human consciousness. Nelson Parker.}.

\section{Les arguments liés à la gouvernance et à la gestion}

La troisième famille de critiques vient des gens dont le métier est de faire fonctionner des organisations — et leur première question est désarmante de simplicité : la hiérarchie servait à quelque chose ; qu'avez-vous prévu à la place ?

Car la hiérarchie, rappellent-ils, n'est pas qu'un instrument de domination. C'est aussi une technologie de coordination : elle clarifie les responsabilités, accélère certaines décisions, offre un recours en cas de litige\footnote{Foss, N. J., \& Klein, P. G. (2014). Why Managers Still Matter. MIT Sloan Management Review, 56(1), 73-80.}. Les cadres intermédiaires, cibles favorites des libérateurs, jouent souvent un rôle décisif que l'on ne découvre qu'après les avoir supprimés — traducteurs entre la stratégie et le terrain, amortisseurs de crises, artisans du changement\footnote{Huy, Q. N. (2001). In praise of middle managers. Harvard Business Review, 79(8), 72-79.}. On peut remplacer ces fonctions. On ne peut pas juste les effacer.

Vient ensuite l'argument de la conduite du changement : la transformation est longue, coûteuse, sans garantie — et la littérature sur le changement organisationnel documente depuis longtemps la fréquence des échecs, même pour des transformations bien moins radicales\footnote{Appelbaum, S. H., Habashy, S., Malo, J. L., \& Shafiq, H. (2012). Back to the future: revisiting Kotter's 1996 change model. Journal of Management Development, 31(8), 764-782.}. Et celui de la réversibilité : que se passe-t-il en cas de crise grave ? Si la réponse est « on recentralise », que reste-t-il de la confiance le lendemain ?

Enfin, la critique qui devrait empêcher tout dirigeant libérateur de dormir : la protection des personnes. Harcèlement, discrimination, inéquité salariale — dans une organisation classique, ces sujets ont des circuits de traitement, imparfaits mais identifiés. Dans une entreprise libérée qui a démonté ses process RH sans les reconstruire, vers qui se tourne une victime ? Le droit du travail, lui, n'a pas été « libéré » : l'employeur reste responsable. Une entreprise libérée sans dispositifs de protection explicites n'est pas audacieuse. Elle est en faute.

Voilà le réquisitoire, présenté sans caricature. Vous remarquerez une chose : presque aucune de ces critiques ne dit « la liberté est une mauvaise idée ». Elles disent « la liberté sans structure, sans garde-fous et sans lucidité fait des dégâts ». Sur ce point précis — et ça m'a coûté de l'admettre — les critiques et moi sommes d'accord.
